\chapter*{Scope and Conventions}
\addcontentsline{toc}{chapter}{Scope and Conventions}

This volume specifies Adversaria in five parts. Part~I states the problem and the treatment of provenance, and Part~II places the design among its predecessors. Part~III, Chapters 4 to 12, is the formal model: the claim and its scope algebra, evidence and warrant, arguments, objections and defeat semantics, distinction and dependency with the gluing analysis, the append-only ledger, the typed relation algebra, constraint, and the status vector. Part~IV, Chapters 13 to 18, is the protocol: it separates routing from adjudication and specifies the pipeline, ontologies and mappings, capability gates and restrictions, imports and negative results, and federation and queries. Part~V works two cases. The parts after the formal model add no semantic rule to it; each chapter states which of its results are proved in Part~III and which are design requirements.

\paragraph{Conventions.} Statements of three kinds are kept apart throughout. A \emph{theorem}, \emph{proposition} or \emph{lemma} is a mathematical fact about the model and is proved. A \emph{definition} or \emph{design rule} is a choice; it is argued for but not proved, and no theorem is claimed for it. A \emph{counterexample} shows why a tempting alternative is rejected. Nothing in the text is presented as a theorem when it is a normative choice about what counts as good evidence.

\paragraph{Illustrative material.} Examples are schematic. They exercise the formalism and do not report empirical findings about any real study, population or page.
